\section{Εισαγωγή}
\label{Introduction}
\subsection{Ορισμός του Προβλήματος}

Η συγκεκριμένη εργασία αντιμετωπίζει το πρόβλημα του ``διαβάσματος'' των χειλιών (lip reading) ή, αλλιώς, της αναγνώρισης της ομιλίας αξιοποιώντας οπτικά στοιχεία, παραδείγματος χάριν από βίντεο ενός ομιλητή, του οποίου ο ήχος είναι επηρεασμένος από θόρυβο ή απουσιάζει πλήρως. Το πρόβλημα αυτό αποτελεί παραλλαγή του κλασσικού προβλήματος της Αυτόματης Αναγνώρισης Ομιλίας (Automatic Speech Recognition - ASR), στόχος του οποίου είναι η μετατροπή της ομιλούμενης γλώσσας σε γραπτή μορφή, αναλύοντας κατάλληλα την κυματομορφή του ηχητικού σήματος. Για τον λόγο αυτό, συχνά δίνεται στο πρόβλημα που αντιμετωπίζεται στη μελέτη αυτή ο τίτλος ``Οπτική Αναγνώριση Ομιλίας'' (Visual Speech Recognition - VSR) ή, σε περίπτωση που συνδυάζονται οπτικά αλλά και ακουστικά στοιχεία ως είσοδοι, ``Οπτικο-Ακουστική Αναγνώριση Ομιλίας'' (Audio-Visual Speech Recognition - AVSR).

Λύση του προβλήματος αποτελεί κατάλληλη μέθοδος η οποία, δεδομένου ενός βίντεο ομιλητή, από το οποίο είτε απουσιάζει (VSR) είτε συνυπάρχει (AVSR) το ηχητικό σήμα, αναγνωρίζει την ομιλία που περιέχεται σε αυτό και την αναπαράγει σε γραπτή μορφή - προβλέπει, δηλαδή, λέξεις και φράσεις που εκφέρει ένας ομιλητής αξιοποιώντας είτε αποκλειστικά οπτικά είτε οπτικοακουστικά στοιχεία. Ιδιαίτερα στην περίπτωση που ο ήχος απουσιάζει πλήρως (VSR), το πρόβλημα αυτό είναι σημαντικά δυσκολότερο από το παραδοσιακό πρόβλημα του ASR, καθώς το σήμα του βίντεο συχνά δεν παρέχει αρκετές φωνολογικές πληροφορίες ώστε η πρόβλεψη λέξης αξιοποιώντας αποκλειστικά αυτό να είναι μονοσήμαντη. Για παράδειγμα, οι αγγλικές λέξεις ``bin'' και ``pin'' παρουσιάζουν πανομοιότυπη όψη κατά την εκφορά τους, με αποτέλεσμα η διαφοροποίηση μεταξύ τους χρησιμοποιώντας αποκλειστικά οπτικά στοιχεία να είναι αδύνατη, ενώ η διαφοροποίηση με ακουστικά μέσα είναι μεν δύσκολη αλλά εφικτή. Το πρόβλημα αυτό είναι κεντρικό στον τομέα του VSR και παρουσιάζεται εκτενώς σε επόμενα κεφάλαια. Επίσης, το βίντεο είναι εν γένει μέσο το οποίο αποτελείται από πολλά περισσότερα δεδομένα από αυτά που είναι πραγματικά χρήσιμα για την αναγνώριση της ομιλίας (π.χ. το παρασκήνιο ενός ομιλητή, ή frames στα οποία ο ομιλητής δεν ομιλεί), σε αντίθεση με τον ήχο, στον οποίο, θεωρώντας έλλειψη θορύβου, ολόκληρο το σήμα περιέχει χρήσιμα φωνολογικά χαρακτηριστικά, όσο ο εκφωνητής ομιλεί.

Στην περίπτωση του προβλήματος του AVSR, έγκυρη λύση αποτελεί οποιαδήποτε μέθοδος που ``παντρεύει'' κατάλληλα την πληροφορία του ηχητικού σήματος με αυτή του βίντεο, ώστε να προβλέπεται ορθά η ομιλία. Σε περιβάλλοντα στα οποία εφαρμόζεται και υπάρχει θόρυβος, αυτό σημαίνει ότι η μέθοδος αναγνωρίζει την αξιοπιστία των δύο modalities και με βάση αυτήν τα συνδυάζει κατάλληλα ώστε να προβλέψει τα λεγόμενα του ομιλητή. Για παράδειγμα, όταν υπάρχουν τεχνικά προβλήματα κατά την καταγραφή του βίντεο (θόλωμα, πάγωμα, ...) ή όταν παρεμβάλλεται εμπόδιο ανάμεσα στην κάμερα και στον ομιλητή, το μοντέλο πρέπει να συμπεραίνει ότι τα χαρακτηριστικά που εξάγονται από το βίντεο είναι μη αξιόπιστα και άρα να βασίζει την πρόβλεψή του κυρίως στο ηχητικό σήμα. Αντίθετα, όταν το ηχητικό σήμα αλλοιώνεται από θόρυβο, λόγω π.χ. βαβούρας μεταξύ των θεατών μίας ομιλίας, ή όταν υπάρχουν τεχνικά προβλήματα κατά την καταγραφή του ήχου, όπως ``κόψιμο'' του μικροφώνου, το μοντέλο οφείλει να αξιοποιεί τα χαρακτηριστικά του βίντεο σε υψηλότερο βαθμό.

Πιθανές εφαρμογές ενός επιτυχημένου μοντέλου που επιλύει το πρόβλημα του VSR ή του AVSR υπάρχουν σε οποιονδήποτε τομέα της ζωής που είναι αναγκαία η επικοινωνία μεταξύ ανθρώπων. Για παράδειγμα, ένα κατάλληλο σύστημα μπορεί να υποστηρίξει την επικοινωνία ατόμων με πλήρη ή μερική αναπηρία, όπως των κωφών ή των βαρήκοων, με τους συνανθρώπους τους, μια και θα μπορούσαν να καταλάβουν την ομιλία δίχως την ανάγκη για νοηματική γλώσσα (την οποία οι περισσότεροι - δυστυχώς - δε γνωρίζουμε). Μία άλλη πιθανή εφαρμογή είναι σε συστήματα Augmented Reality (AR), ώστε να ενισχύσει την κατανόηση της ομιλίας σε καθημερινά περιβάλλοντα, τα οποία συχνά χαρακτηρίζονται από θόρυβο. Σκεφτείτε ένα αμφιθέατρο με φοιτητές που προκαλούν τη συνηθισμένη βαβούρα - το σύστημα θα αξιοποιούσε τόσο τα οπτικά όσο και τα ακουστικά στοιχεία της ομιλίας του καθηγητή ή της καθηγήτριας, επιτρέποντας στους φοιτητές να αντιλαμβάνονται το περιεχόμενο του λόγου του/της. Ακόμη και σε τομείς που ο θόρυβος δεν οφείλεται σε ανθρώπινη ομιλία αλλά μηχανήματα, π.χ. εργοτάξια, μπορεί να αξιοποιηθεί ένα σύστημα AVSR για να ενισχύεται η ικανότητα των ανθρώπων να επικοινωνούν και άρα να επιταχύνεται και να γίνεται πιο ασφαλές το έργο τους.

\subsection{Διαθέσιμες Μέθοδοι}

Τo VSR ή, όπως αλλιώς ονομάζεται, Automatic Lip Reading, έχει γνωρίσει τα τελευταία χρόνια σημαντική πρόοδο, ξεκινώντας από απλά μοντέλα ταξινόμησης απομονωμένων frames σε φωνήματα ή visemes (βλ. \ref{Phonemes}) και φτάνοντας πλέον σε συστήματα end-to-end που αξιοποιούν τη βαθιά μάθηση (deep learning) και προβλέπουν από το βίντεο απ'ευθείας το περιεχόμενο της ομιλίας. Αρχιτεκτονικές όπως το LipNet \cite{assael2016lipnetendtoendsentencelevellipreading} και το Watch, Listen, Attend and Spell (WLAS, \cite{Chung17} καθιέρωσαν τη χρήση 3D συνελικτικών νευρωνικών δικτύων (τόσο στους 2 χωρικούς όσο και στον χρονικό άξονα), ακολουθούμενα από decoders βασισμένους είτε στη λογική των recurrent νευρωνικών δικτύων είτε στον μηχανισμό του attention, εκπαιδευόμενων είτε με Connectionist Temporal Classification (βλ. \ref{MachineLearningBackground}) είτε με αντικειμενικές συναρτήσεις sequence-to-sequence. Μετέπειτα δουλειές έχουν αξιοποιήσει τους Transformers ως temporal encoders, ενώ άλλες σημαντικές μελέτες έχουν στραφεί προς την προεκπαίδευση οπτικο-ακουστικών μοντέλων σε μεγάλη κλίμακα με self-supervised προβλήματα, όπως το AV-HuBERT \cite{shi2022learningaudiovisualspeechrepresentation}, το οποίο μαθαίνει κοινές αναπαραστάσεις του ήχου και του βίντεο από unlabelled δεδομένα, ή την αυτόματη δημιουργία labels και την ενίσχυση της επίβλεψης, όπως το Auto-AVSR \cite{Ma_2023}, που χρησιμοποιεί ήδη υπάρχοντα ASR μοντέλα για την επέκταση των labelled δεδομένων.

Πρόσφατα έχει στραφεί η προσοχή στον διαχωρισμό του βασικού προβλήματος του VSR σε δύο υπο-προβλήματα, αφενός την πρόβλεψη των φωνημάτων (βλ. \ref{Phonemes}) και αφετέρου τον συνδυασμό των προβλεπόμενων φωνημάτων για την ανακατασκευή λέξεων και, τελικά, φράσεων. Τέτοιο είναι και το σύστημα που προτείνεται από το VALLR \cite{thomas2026vallrvisualasrlanguage}, το οποίο πετυχαίνει state-of-the-art επιδόσεις σε πρόβλεψη λέξεων, χρησιμοποιώντας πολύ λιγότερα labelled δεδομένα σε σχέση με τα μοντέλα με τις αμέσως καλύτερες επιδόσεις. Σε αυτήν τη μεθοδολογία βασίζεται σε μεγάλο βαθμό η δουλειά της συγκεκριμένης διπλωματικής. Όπως περιγράφεται και παρακάτω στην ενότητα \ref{ThisWork}, η εργασία αυτή εστιάζει κυρίως στο υπο-πρόβλημα της πρόβλεψης φωνημάτων και όχι τόσο στο δεύτερο υπο-πρόβλημα, αυτό του συνδυασμού των φωνημάτων. 

Εκτενέστερη αναφορά στο σχετικό State of the Art για το πρόβλημα του VSR και του AVSR βρίσκεται στο κεφάλαιο \ref{StateOfTheArt}.

\subsection{Περιγραφή της Συγκεκριμένης Εργασίας}
\label{ThisWork}

Όπως αναφέρθηκε και παραπάνω, η εργασία αυτή εστιάζει σε σημαντικό βαθμό στην επίλυση του πρώτου υποπροβλήματος, την πρόβλεψη, δηλαδή, φωνημάτων είτε αποκλειστικά από το βίντεο, μέσω του visual μοντέλου (βλ. \ref{ModelStructure}), είτε μέσω του βίντεο και του ήχου, αξιοποιώντας το audiovisual μοντέλο (βλ. \ref{AudiovisualModel}). Για την επίλυση του 2ου προβλήματος, του συνδυασμού των προβλεπόμενων φωνημάτων ώστε να σχηματιστούν λέξεις και φράσεις, δε δαπανήθηκε ιδιαίτερος χρόνος, καθώς χρησιμοποιήθηκε αποκλειστικά για να παρουσιαστούν τα αποτελέσματα των μοντέλων με μορφή πιο κατανοητή και όχι για να δοκιμαστεί κάποια νέα μέθοδος.

Αναπτύχθηκε, αρχικά, το visual μοντέλο, το οποίο αποτελείται από ένα Front-End συνελικτικού δικτύου τριών διαστάσεων, ακολουθούμενο από BiLSTM Sequence Encoder και μίας γραμμικής στρώσης ταξινόμησης. Το μοντέλο εκπαιδεύεται μέσω αντικειμενικής συνάρτησης που περιλαμβάνει επίβλεψη τόσο σε επίπεδο frame, μέσω ενός όρου cross-entropy, όσο και σε επίπεδο αλληλουχίας φωνημάτων, αξιοποιώντας έναν όρο CTC. Το μοντέλο παρουσιάζεται αναλυτικότερα στην ενότητα \ref{ModelStructure}. 

Για την εκπαίδευση του μοντέλου αυτού, καθώς χρησιμοποιείται όρος cross-entropy, ήταν αναγκαία η γνώση ποιου φωνήματος εκφωνείται σε κάθε frame του κάθε βίντεο. Αναπτύχθηκε αλγόριθμος εξαγωγής φωνημάτων για κάθε frame, αξιοποιώντας τόσο την πραγματική φράση που εκφωνείται (εφ' όσον είναι γνωστή) όσο και pre-trained μοντέλα που προβλέπουν φωνήματα από τον ήχο του κάθε βίντεο. Έτσι, συνδυάζεται η πληροφορία των αναμενόμενων φωνημάτων από την ground truth φράση και του χρονισμού εκφοράς από την πρόβλεψη του pre-trained μοντέλου. Για μείωση των δεδομένων που εισέρχονται στο μοντέλο, εξήχθη από κάθε βίντεο η περιοχή των χειλιών, χρησιμοποιώντας πάλι pre-trained μοντέλα. Αναπτύχθηκε μέθοδος δημιουργίας soft targets για τα φωνήματα, με βάση την ομοιότητά τους (οπτική και ακουστική). Εφαρμόστηκε data augmentation στα βίντεο ώστε να αυξηθεί η ευρωστία της εκπαίδευσης. Τα βήματα αυτά εξηγούνται αναλυτικότερα στην ενότητα \ref{DataPreprocessing}. Αναπτύχθηκε αλγόριθμος πρόβλεψης φράσεων από αλληλουχία φωνημάτων, δεδομένων των περιορισμών στο λεξιλόγιο και στο συντακτικό ενός dataset, ο οποίος παρουσιάζεται στην υποενότητα \ref{PostProcessing}.

Μετά, αναπτύχθηκε το audiovisual μοντέλο, με λογική παρόμοια με αυτή του visual μοντέλου. Αποτελείται από κατάλληλα συνελικτικά νευρωνικά δίκτυα για την πληροφορία του βίντεο και του ήχου. Έπειτα, εφαρμόζεται μηχανισμός Attention μεταξύ των features εξαγόμενων από το βίντεο και από τον ήχο, ώστε να συνδυάζονται οι πληροφορίες αυτές. Ακολουθεί συνένωση των δύο feature vectors με κατάλληλα gates που επηρεάζονται από την αξιοπιστία του κάθε modality κάθε χρονική στιγμή (πληροφορία που θεωρείται γνωστή, άρα δε χρειάζεται να προβλεφθεί από το μοντέλο). Τέλος, ακολουθεί Sequence Encoder (είτε BiLSTM είτε Transformer) και γραμμικός ταξινομητής. Το μοντέλο και πάλι εκπαιδεύεται χρησιμοποιώντας συνάρτηση με όρους cross-entropy και CTC. Το μοντέλο παρουσιάζεται αναλυτικότερα στην ενότητα \ref{AudiovisualModel}. 

Ειδικά για το μοντέλο αυτό, χρησιμοποιούνται μερικά περαιτέρω βήματα Pre-Processing, συγκεκριμένα η πρόβλεψη της αξιοπιστίας του βίντεο και του ήχου για κάθε χρονική στιγμή, η προσθήκη θορύβου σε οποιοδήποτε από τα δύο modalities και ο υπολογισμός των Mel-Spectrograms των ηχητικών σημάτων, τα οποία αποτελούν την είσοδο της πληροφορίας του ήχου στο μοντέλο. Τα βήματα αυτά παρουσιάζονται στην υποενότητα \ref{AudiovisualPreprocessing}.

Η επίδοση των μοντέλων αυτών δοκιμάστηκε σε μία σειρά 3 πειραμάτων. Το πρώτο πείραμα (βλ. \ref{Experiment1}) αποτελείται από 2 υπο-πειράματα, το πρώτο εκ των οποίων αφορά μία μελέτη αφαίρεσης (ablation study) της προτεινόμενης μεθοδολογίας για το visual μοντέλο. Συγκεκριμένα, δοκιμάζονται μοντέλα με διαφορετικό Sequence Encoder (BiLSTM, Transformer ή καθόλου), με διαφορετικές loss functions (μόνο CTC, μόνο cross-entropy ή συνδυασμός των δύο) και με χρήση ή χωρίς των soft targets. Το μοντέλο που αναδεικνύεται βέλτιστο από αυτή τη διαδικασία δοκιμάζεται περαιτέρω στο υπο-πείραμα 2, κατά το οποίο εφαρμόζεται και η πρόβλεψη λέξεων (πέραν των φωνημάτων). Χρησιμοποείται το GRID dataset (βλ. \ref{Datasets}), περιβάλλον δηλαδή περιορισμένο ως προς το λεξιλόγιο και το συντακτικό.

Το δεύτερο πείραμα (βλ. \ref{Experiment2}) είχε σκοπό τη δοκιμή του προτεινόμενου visual μοντέλου σε in the wild περιβάλλον, με τη χρήση του ευρέως διαδεδομένου στον τομέα LRS2 dataset. Παρά το γεγονός ότι ακολουθήθηκαν τα απαιτούμενα βήματα και η συμπλήρωση αίτησης για απόκτηση πρόσβασης στο συγκεκριμένο dataset, έως τώρα δεν έχει ληφθεί επιβεβαίωση για την άδεια χρήσης και ως εκ τούτου ήταν αδύνατον να ολοκληρωθεί το πείραμα. Συμπεριλήφθηκαν, για λόγους πληρότητας, οι επιδόσεις άλλων state of the art μοντέλων, ενώ αναφέρθηκε η διεξαγωγή του συγκεκριμένου πειράματος ως η βασική μελλοντική επέκταση.

Τέλος, το τρίτο πείραμα (βλ. \ref{Experiment3}) αφορούσε τη δοκιμή του audiovisual μοντέλου σε περιβάλλον περιορισμένο (χρησιμοποιήθηκε και πάλι το GRID dataset, βλ. \ref{Datasets}), αλλά υπό καθεστώς ύπαρξης θορύβου τόσο στο βίντεο όσο και στον ήχο. Συγκεκριμένα, στο πρώτο υπο-πείραμα μελετήθηκε η επίδραση της αύξησης του SNR του ηχητικού σήματος στην επίδοση του μοντέλου, ενώ στο δεύτερο υπο-πείραμα παρουσιάστηκαν τα αποτελέσματα της μείωσης της αξιοπιστίας του βίντεο στην ακρίβεια των προβλέψεων, σε επίπεδο φωνημάτων.

\subsection{Δομή της Συγκεκριμένης Αναφοράς}

Παρακάτω περιγράφεται εν συντομία η δομή της συγκεκριμένης αναφοράς της διπλωματικής εργασίας.

Προηγήθηκαν στην ενότητα \ref{Summaries} οι περιλήψεις της εργασίας στα ελληνικά και στα αγγλικά, ενώ στο συγκεκριμένο κεφάλαιο (\ref{Introduction}) έγινε μία εισαγωγή, ορίζοντας το πρόβλημα του Lip Reading, παρουσιάζοντας μερικές από τις διαθέσιμες μεθόδους στη βιβλιογραφία και περιγράφοντας σύντομα το αντικείμενο μελέτης της εργασίας.

Ακολουθεί το κεφάλαιο με τίτλο ``Τεχνικό Υπόβαθρο'' (\ref{Background}), στην οποία εξηγούνται οι έννοιες αναγκαίες για την κατανόηση της εργασίας. Γίνεται, αρχικά, μία αναφορά στα φωνήματα και στα visemes (βλ. \ref{Phonemes}), η οποία ακολουθείται από μία σύντομη παρουσίαση μερικών από των εννοιών της μηχανικής μάθησης, συγκεκριμένα του Positional Encoding, της συνάρτησης κόστους Cross-Entropy και της συνάρτησης κόστους CTC (βλ. \ref{MachineLearningBackground}). Να σημειωθεί ότι οι συναρτήσεις κόστους παρουσιάζονται κυρίως ώστε να αποφευχθεί σύγχυση για τον τρόπο με τον οποίον υπολογίζονται στα πειράματα, ενώ δεν παρουσιάζονται σε αυτήν την ενότητα έννοιες της μηχανικής μάθησης πιο γενικές, όπως τα CNNs, τα BiLSTM, οι Transformers, το Dropout, το Layer Normalization και άλλες, για να μειωθεί η έκταση του κειμένου, καθώς θεωρήθηκαν γνωστές. Τέλος, παρουσιάζονται και 2 ακόμη έννοιες, η απόσταση Levenshtein και τα Mel-Spectrograms (βλ. \ref{OtherBackground}), τα οποία χρησιμοποιούνται στην εργασία και δεν ταιριάζουν σε κάποια από τις προαναφερθείσες κατηγορίες.

Έπειτα, ακολουθεί μία εκτενέστερη αναφορά στο State of the Art του προβλήματος του Lip Reading από αυτήν που γίνεται στο συγκεκριμένο κεφάλαιο, στο κεφάλαιο \ref{StateOfTheArt} με τίτλο State of the Art.

Το επόμενο κεφάλαιο είναι το μεγαλύτερο σε μέγεθος και έχει τίτλο Μεθοδολογία \ref{Methodology}. Παρουσιάζονται αναλυτικά όλα τα βήματα της μεθοδολογίας που εφαρμόζεται και δοκιμάζεται, τόσο για το visual όσο και για το audiovisual μοντέλο, αρχίζοντας από την προεπεξεργασία των δεδομένων (βλ. \ref{DataPreprocessing}). Παρουσιάζεται η αρχιτεκτονική πρώτα του visual μοντέλου (βλ. \ref{ModelStructure}) και έπειτα του audiovisual μοντέλου (βλ. \ref{AudiovisualModel}), με κατάλληλο σχολιασμό των σχεδιαστικών επιλογών. Τέλος, γίνεται αναφορά στο post-processing των δεδομένων στην ενότητα \ref{PostProcessing}.

Στο κεφάλαιο με τίτλο ``Αξιολόγηση της Μεθοδολογίας'' (\ref{Evaluation}), παρουσιάζονται αναλυτικά τα διάφορα πειράματα που διεξήχθησαν για τη συγκεκριμένη εργασία, αφού πρώτα παρουσιαστούν τα Datasets (βλ. \ref{Datasets}) και τα Κριτήρια Αξιολόγησης (βλ. \ref{Criteria}). Στις υποενότητες \ref{Experiment1}, \ref{Experiment2} και \ref{Experiment 3} παρουσιάζονται με τη σειρά τα 3 πειράματα που διεξήχθησαν.

Τέλος, στο κεφάλαιο με τίτλο ``Συμπεράσματα και Μελλοντικές Επεκτάσεις'' (\ref{ConclusionsFuture}), παρουσιάζονται τα συμπεράσματα (βλ. \ref{Conclusions}) και οι μελλοντικές επεκτάσεις (βλ. \ref{Future}) που προκύπτουν από την εργασία.
